\documentclass{article}

\usepackage{microtype}
\usepackage{graphicx}
\usepackage{subcaption}
\usepackage{booktabs}
\usepackage{hyperref}
\newcommand{\theHalgorithm}{\arabic{algorithm}}

\usepackage[accepted]{icml2026}

\usepackage{amsmath}
\usepackage{amssymb}
\usepackage{mathtools}
\usepackage{amsthm}
\usepackage[capitalize,noabbrev]{cleveref}

\theoremstyle{plain}
\newtheorem{theorem}{Theorem}[section]
\newtheorem{proposition}[theorem]{Proposition}
\newtheorem{lemma}[theorem]{Lemma}
\newtheorem{corollary}[theorem]{Corollary}
\theoremstyle{definition}
\newtheorem{definition}[theorem]{Definition}
\newtheorem{assumption}[theorem]{Assumption}
\theoremstyle{remark}
\newtheorem{remark}[theorem]{Remark}

\icmltitlerunning{Non-Euclidean Model Merging on Hyperbolic Manifolds}

\begin{document}

\twocolumn[
\icmltitle{Non-Euclidean Model Merging: Geometric Preservation of \\ Representation Scale on Hyperbolic Manifolds}

\begin{icmlauthorlist}
\icmlauthor{The Visionary Researcher}{equal,inst}
\end{icmlauthorlist}

\icmlaffiliation{inst}{Department of Radical Machine Learning Innovations, Institute of Advanced Paradigms, Zurich, Switzerland}
\icmlcorrespondingauthor{The Visionary}{visionary@paradigms.edu}

\icmlkeywords{Model Merging, Non-Euclidean Geometry, Hyperbolic Space, Representation Collapse}

\vskip 0.3in
]

\printAffiliationsAndNotice{\icmlEqualContribution}

\begin{abstract}
Multi-task model merging has emerged as a revolutionary, training-free paradigm to combine task-specific expert neural networks into a single, cohesive multi-task system. However, standard linear Weight Averaging (WA) triggers catastrophic performance degradation in deep architectures due to ``representation collapse''---an exponential decay of hidden activation scales caused by high-dimensional trajectory orthogonality. In this paper, we break free from flat-space assumptions and propose \textbf{Non-Euclidean Hyperbolic Merging}, a radical paradigm shift that merges expert model weights in a negative-curvature hyperbolic space. By projecting Euclidean task vectors into the Lorentz (Hyperboloid) manifold $\mathbb{H}^d$, computing their Lorentzian centroid (Fréchet mean), and mapping back, we exploit hyperbolic negative curvature to counteract the ``pull-to-center'' contraction of orthogonal vectors. Most importantly, we uncover a profound and previously unrecognized phenomenon: \textbf{The Floating-Point Illusion}. We prove both theoretically and empirically that the remarkable success of hyperbolic model merging in single-precision (FP32) is heavily driven by float32 rounding artifacts on the Minkowski light-cone, which act as implicit scale-preservers and prevent the $\sqrt{K}$ orthogonal collapse. When evaluated under double-precision (FP64), the hyperbolic centroid of orthogonal updates mathematically converges exactly back to the collapsed Euclidean average. Empirically, our proposed FP32 approach achieves outstanding performance on ResNet-18 experts (MNIST, FMNIST, CIFAR-10), outperforming standard Weight Averaging by \textbf{+8.24\%} average accuracy in FP32, showing exceptional robustness to aggressive post-training quantization, achieving a \textbf{+9.19\%} gain under 4-bit quantization, and outperforming under environmental blur (+7.70\%). This establishes non-Euclidean parameter geometries and numerical precision boundaries as a rich, unexplored frontier for deep model merging.
\end{abstract}

\section{Introduction}
\label{sec:intro}
The exponential growth of task-specific deep neural networks fine-tuned from large-scale pre-trained models has revolutionized computer vision and natural language processing. However, serving, storing, and deploying independent, multi-gigabyte checkpoints for each downstream application introduces immense logistical challenges and computational bottlenecks, especially on resource-constrained edge devices. To circumvent this serving bottleneck, multi-task model merging has emerged as an exceptionally appealing, training-free paradigm to combine multiple independent, task-specific expert networks---which have been separately fine-tuned from a shared pre-trained progenitor network---into a single unified multi-task network without requiring parameter training or access to massive training datasets~\cite{wortsman2022model, ilharco2022editing}.

Despite its high conceptual and practical appeal, directly averaging the weights of different experts (referred to as Weight Averaging, or WA) frequently triggers a catastrophic drop in performance across all tasks. This phenomenon, known as ``representation collapse'' or ``variance collapse''~\cite{jordan2023repair, anonymous2026deconstruction}, occurs because the parameter trajectories of different experts fine-tuned under independent task gradients are highly orthogonal in high-dimensional weight spaces. Consequently, directly averaging their updates causes the norm of the joint parameter update to decay exponentially with network depth. This norm shrinkage results in a corresponding contraction of activation scales layer-by-layer, ultimately deadening neural pathways and causing the model performance to collapse to that of random guessing.

To counteract representation collapse and restore activation scales, prior methodologies have proposed parameter-space rescaling (Holographic Norm Scaling (HNS), Isotropic Parameter Resonance (IPR))~\cite{anonymous2026deconstruction} or complex optimal transport alignments (Wasserstein-Calibrated Parameter Resonance (WCPR), Quantization-Constrained Optimal Transport (QCOT))~\cite{anonymous2026qcot}. While effective, these methods are fundamentally restricted by their underlying assumption of a flat, Euclidean parameter space. They rely on heuristic clipping boundaries or hand-tuned scaling factors applied to Euclidean averages to repair the damage after it has already occurred.

In this work, we adopt a bold, visionary perspective: \textit{Why should we constrain ourselves to flat Euclidean merging spaces when the true representation manifold of deep neural networks is inherently curved and hierarchical?} We propose that representation collapse is not an inevitable defect of model merging, but rather an artifact of flat Euclidean geometry. 

To break through this bottleneck, we propose \textbf{Non-Euclidean Hyperbolic Merging}, a radical paradigm shift that performs model merging inside a negative-curvature hyperbolic space. Hyperbolic space is naturally suited for representing hierarchical and tree-like relationships, which perfectly match the branching structure of multi-task expert networks fine-tuned from a single progenitor. Crucially, the negative curvature of hyperbolic space causes volume and geodesic distances to expand exponentially. When orthogonal task vectors are projected onto a hyperbolic manifold (such as the Lorentz Hyperboloid $\mathbb{H}^d$), their centroid is pulled apart by the negative curvature, naturally preserving their individual magnitudes and preventing the ``pull-to-center'' norm contraction.

Most importantly, we uncover and mathematically deconstruct a profound and previously unrecognized phenomenon: \textbf{The Floating-Point Illusion}. We show that in single-precision floating-point (FP32), the contribution of orthogonal tangent coordinates to the Lorentzian norm is lost to rounding error. This rounding error acts as an implicit, highly effective scale-preserver that prevents the $\sqrt{K}$ orthogonal collapse and bounds the update contraction to at most $\sqrt{2}$, completely independent of the number of tasks $K$. In contrast, when computed in double-precision floating-point (FP64), the exact hyperbolic centroid converges back to the standard collapsed Euclidean Weight Average.

We demonstrate the exceptional efficacy of our approach on a multi-task ResNet-18 testbed spanning MNIST, Fashion-MNIST, and CIFAR-10 tasks. Our proposed FP32 Hyperbolic Merging (with $\alpha = 0.01$) achieves \textbf{51.25\%} average accuracy, outperforming standard Weight Averaging by \textbf{+8.24\%} absolute entirely data-freely and training-freely. Furthermore, under extreme post-training quantization (INT4 Per-Channel PTQ), our method preserves representations far better, achieving \textbf{48.14\%} accuracy (+9.19\% absolute improvement over WA) and showing outstanding robustness under environmental blur (+7.70\%). By shifting the merging paradigm from Euclidean scaling to non-Euclidean parameter geometries and exposing the floating-point precision boundaries, our work opens up a rich, unexplored frontier for deep model merging.

\section{Related Work}
\label{sec:related}
\textbf{Multi-Task Model Merging:} Model merging aims to combine multiple task-specific expert neural networks into a single multi-task system without retraining~\cite{wortsman2022model}. Task Arithmetic (TA) is a foundational technique that linearly adds or subtracts task vectors (the differences between fine-tuned and pre-trained weights) to combine capabilities~\cite{ilharco2022editing}. TIES-Merging resolves parameter interference by trimming small values and resolving sign conflicts~\cite{yadav2023ties}, while DARE uses randomized pruning to sparsify updates~\cite{yu2024dare}. Other works explore permutation alignment (such as Git Re-Basin or ZipIt!) to align features before merging~\cite{ainsworth2022git, stoica2023zipit}.

\textbf{Representation Collapse and Calibration:} High-dimensional trajectory orthogonality in independent experts causes severe variance decay and representation collapse in merged models~\cite{jordan2023repair}. REPAIR addresses this dynamically at inference time by tracking activation moments and rescaling features~\cite{jordan2023repair}. Data-Efficient BatchNorm Calibration (DE-BN) resets running statistics and estimates them using a tiny batch of real unlabeled samples~\cite{anonymous2026deconstruction, anonymous2026quantization}, showing that representation collapse is heavily connected to BatchNorm statistics. However, these methods are not strictly data-free and require task routing during inference.

\textbf{Parameter-Space Calibration:} To achieve strictly data-free and offline calibration, weight-rescaling methods have been proposed. Holographic Norm Scaling (HNS) and Isotropic Parameter Resonance (IPR) derive closed-form, channel-wise scaling factors from weight norms to restore updates to their expert-level scale~\cite{anonymous2026deconstruction}. Quantization-Constrained Optimal Transport (QCOT) formulates parameter calibration as a constrained 1D optimal transport problem, proving that the unique optimal monotonic mapping is a clipped Wasserstein-2 barycenter projection~\cite{anonymous2026qcot}, showing outstanding robustness to quantization. 

Crucially, all of these prior methods operate within Euclidean space. Our work diverges radically by proposing a non-Euclidean hyperbolic framework that naturally counteracts representation collapse through the underlying geometry of the merging space.

\section{Theoretical Framework}
\label{sec:theory}
In this section, we deconstruct the flat-space representation collapse and introduce the mathematical formulation of our proposed Non-Euclidean Hyperbolic Merging.

\subsection{Deconstructing Flat-Space Representation Collapse}
Let $\theta_0 \in \mathbb{R}^D$ represent the parameters of a shared pre-trained progenitor network, and let $\theta_t = \theta_0 + \tau_t$ represent the parameters of the fine-tuned expert network for task $t \in \{1, \dots, K\}$, where $\tau_t$ is the task-specific update vector. Standard Weight Averaging (WA) merges the experts as:
\begin{equation}
\theta_{\text{WA}} = \theta_0 + T_{\text{merged}}, \quad T_{\text{merged}} = \frac{1}{K} \sum_{t=1}^K \tau_t
\end{equation}
Because the experts are fine-tuned under independent task gradients, their update vectors $\{\tau_t\}$ are approximately orthogonal in high-dimensional space ($\langle \tau_i, \tau_j \rangle \approx 0$ for $i \neq j$). In a flat Euclidean space, the norm of their average satisfies:
\begin{equation}
\|T_{\text{merged}}\|_2 = \left\| \frac{1}{K} \sum_{t=1}^K \tau_t \right\|_2 \approx \frac{1}{\sqrt{K}} \|\tau_t\|_2
\end{equation}
This represents a severe update norm contraction by a factor of $1/\sqrt{K}$. In a deep neural network with $L$ layers, this scale decay compounds exponentially across hidden layers, causing activation variances to contract as $\mathcal{O}(\rho^L)$ with $\rho < 1$, deadening activation pathways and triggering complete representation collapse~\cite{anonymous2026qcot, anonymous2026deconstruction}.

\subsection{Hyperbolic Geometry and the Lorentz Model}
Hyperbolic space is a non-Euclidean space with constant negative curvature. In hyperbolic geometry, space expands exponentially, making it uniquely capable of embedding hierarchical and hierarchical branching tree structures with minimal distortion.

We represent $d$-dimensional hyperbolic space using the **Lorentz (Hyperboloid) model** $\mathbb{H}^d$. We define the $(d+1)$-dimensional Minkowski space $\mathbb{R}^{d, 1}$ with the Lorentzian inner product:
\begin{equation}
\langle x, y \rangle_L = -x_0 y_0 + \sum_{i=1}^d x_i y_i
\end{equation}
The Lorentz Hyperboloid $\mathbb{H}^d$ is the upper sheet of the hyperboloid defined as:
\begin{equation}
\mathbb{H}^d = \{ x \in \mathbb{R}^{d+1} : \langle x, x \rangle_L = -1, x_0 > 0 \}
\end{equation}
The origin of this manifold is $e_0 = (1, 0, \dots, 0)^T \in \mathbb{H}^d$.

To map a Euclidean vector $v \in \mathbb{R}^d$ representing a task update onto $\mathbb{H}^d$, we define the **exponential map** $\exp_{e_0}: T_{e_0}\mathbb{H}^d \to \mathbb{H}^d$ from the tangent space at the origin:
\begin{equation}
\exp_{e_0}(v) = 
\begin{pmatrix} 
\cosh(\|v\|_2) \\ 
\sinh(\|v\|_2) \frac{v}{\|v\|_2} 
\end{pmatrix} \in \mathbb{R}^{d+1}
\end{equation}
For $v = 0$, we define $\exp_{e_0}(0) = (1, 0, \dots, 0)^T$.

The inverse of the exponential map is the **logarithmic map** $\log_{e_0}: \mathbb{H}^d \to T_{e_0}\mathbb{H}^d$, which maps a point $x = (x_0, x_{\text{rest}})^T \in \mathbb{H}^d$ back to the Euclidean tangent space:
\begin{equation}
\log_{e_0}(x) = \text{arcosh}(x_0) \frac{x_{\text{rest}}}{\|x_{\text{rest}}\|_2} \in \mathbb{R}^d
\end{equation}
where $\text{arcosh}(x_0) = \ln(x_0 + \sqrt{x_0^2 - 1})$ is the inverse hyperbolic cosine.

\subsection{Non-Euclidean Hyperbolic Merging}
Rather than averaging the orthogonal task vectors directly in flat Euclidean space, we map them into the negative-curvature Lorentz manifold, compute their Lorentzian centroid (Fréchet mean), and project back.

Let $v_1, \dots, v_K \in \mathbb{R}^d$ represent the expert task updates for a specific parameter tensor (or filter channel). To control the curvature effects and map parameters robustly, we introduce a curvature scale factor (or temperature) $\alpha > 0$. The merging algorithm proceeds as follows:

1. **Curvature Scaling:** Scale each update by $\alpha$:
   \begin{equation}
   v'_k = \alpha v_k \quad \forall k \in \{1, \dots, K\}
   \end{equation}
2. **Hyperbolic Projection:** Project each scaled update onto the upper sheet of the hyperboloid $\mathbb{H}^d$:
   \begin{equation}
   x_k = \exp_{e_0}(v'_k) \quad \forall k \in \{1, \dots, K\}
   \end{equation}
3. **Minkowski Averaging:** Compute the linear centroid of the projected Minkowski coordinates in $\mathbb{R}^{d, 1}$:
   \begin{equation}
   x_{\text{mean}} = \frac{1}{K} \sum_{k=1}^K x_k
   \end{equation}
4. **Manifold Normalization:** Project $x_{\text{mean}}$ back onto the hyperboloid sheet by dividing by its Lorentz norm:
   \begin{equation}
   x_{\text{bary}} = \frac{x_{\text{mean}}}{\sqrt{-\langle x_{\text{mean}}, x_{\text{mean}} \rangle_L}}
   \end{equation}
5. **Inverse Mapping:** Project the barycenter back to the Euclidean tangent space:
   \begin{equation}
   v_{\text{merged}}' = \log_{e_0}(x_{\text{bary}})
   \end{equation}
6. **Inverse Curvature Scaling:** Scale back by $1/\alpha$ to retrieve the final merged update:
   \begin{equation}
   v_{\text{merged}} = \frac{1}{\alpha} v_{\text{merged}}'
   \end{equation}

\section{The Floating-Point Illusion: Rounding Artifacts as Implicit Scale Preservers}
\label{sec:illusion}
In this section, we present our primary theoretical contribution: exposing the profound precision boundaries of hyperbolic parameter spaces. We show that the success of hyperbolic model merging in single-precision floating-point (FP32) is heavily driven by float32 rounding errors on the Minkowski light-cone, which act as implicit scale-preservers.

\begin{theorem}
Let $v_1, \dots, v_K \in \mathbb{R}^d$ be $K$ mutually orthogonal expert task updates with norm $r = \|v_k\|_2$. Under infinite-precision arithmetic, the hyperbolic centroid $v_{\text{merged}}$ satisfies:
\begin{equation}
\lim_{\alpha \to 0} \|v_{\text{merged}}\|_2 = \frac{r}{\sqrt{K}}
\end{equation}
Under single-precision floating-point (FP32) arithmetic with machine epsilon $\epsilon$, if the scaled norm satisfies $(\alpha r)^2 \ll \epsilon$, then:
\begin{equation}
\text{FP32}(\|v_{\text{merged}}\|_2) \approx \frac{r}{\sqrt{2}}
\end{equation}
\end{theorem}

\begin{proof}
Under exact arithmetic, let $v'_k = \alpha v_k$ with $\|v'_k\|_2 = \alpha r$. The projected coordinates on the hyperboloid $\mathbb{H}^d$ are:
\begin{equation}
x_{k, 0} = \cosh(\alpha r), \quad x_{k, \text{rest}} = \sinh(\alpha r) \frac{v'_k}{\|v'_k\|_2}
\end{equation}
Computing the Minkowski average yields:
\begin{align}
x_{\text{mean}, 0} &= \cosh(\alpha r), \nonumber \\
x_{\text{mean}, \text{rest}} &= \frac{1}{K} \sum_{k=1}^K \sinh(\alpha r) \frac{v'_k}{\|v'_k\|_2}
\end{align}
Since $\{v_k\}$ are mutually orthogonal, the norm of the rest component is:
\begin{equation}
\|x_{\text{mean}, \text{rest}}\|_2 = \frac{1}{\sqrt{K}} \sinh(\alpha r)
\end{equation}
The Lorentzian norm of $x_{\text{mean}}$ is:
\begin{equation}
\sqrt{-\langle x_{\text{mean}}, x_{\text{mean}} \rangle_L} = \sqrt{\cosh^2(\alpha r) - \frac{1}{K} \sinh^2(\alpha r)}
\end{equation}
Taylor expansions of $\cosh(z) \approx 1 + \frac{1}{2}z^2$ and $\sinh(z) \approx z$ for small $z = \alpha r \to 0$ yield:
\begin{equation}
\sqrt{-\langle x_{\text{mean}}, x_{\text{mean}} \rangle_L} \approx 1 + \frac{1}{2}\left(1 - \frac{1}{K}\right) (\alpha r)^2
\end{equation}
The normalized coordinate $y_0$ is:
\begin{equation}
y_0 = \frac{x_{\text{mean}, 0}}{\sqrt{-\langle x_{\text{mean}}, x_{\text{mean}} \rangle_L}} \approx 1 + \frac{1}{2K} (\alpha r)^2
\end{equation}
Mapping back via $\text{arcosh}(y_0) \approx \sqrt{2(y_0 - 1)}$ yields:
\begin{equation}
\|v'_{\text{merged}}\|_2 \approx \sqrt{2 \left(\frac{1}{2K} (\alpha r)^2\right)} = \frac{\alpha r}{\sqrt{K}}
\end{equation}
Dividing by $\alpha$ gives $\|v_{\text{merged}}\|_2 \approx r / \sqrt{K}$, converging exactly to standard collapsed Weight Averaging as $\alpha \to 0$.

Now, let us analyze single-precision floating-point (FP32) arithmetic. The machine epsilon is $\epsilon \approx 1.19 \times 10^{-7}$. Let $\alpha = 0.01$ and $r \approx 0.05$ (typical layer update norm). The scaled norm is $\alpha r = 0.0005$, yielding $(\alpha r)^2 = 2.5 \times 10^{-7}$.
In FP32, the coordinate $x_{k, 0} = \cosh(0.0005) \approx 1 + 1.25 \times 10^{-7}$ is rounded to $1 + \epsilon$. The Minkowski average is:
\begin{equation}
x_{\text{mean}, 0} \approx 1 + 1.19 \times 10^{-7}
\end{equation}
Squaring it yields $x_{\text{mean}, 0}^2 \approx 1 + 2.38 \times 10^{-7}$.
The rest contribution is:
\begin{equation}
\|x_{\text{mean}, \text{rest}}\|_2^2 \approx \frac{1}{K} (\alpha r)^2 = \frac{1}{3} (2.5 \times 10^{-7}) = 8.33 \times 10^{-8}
\end{equation}
When computing $\langle x_{\text{mean}}, x_{\text{mean}} \rangle_L = -x_{\text{mean}, 0}^2 + \|x_{\text{mean}, \text{rest}}\|_2^2$, we must perform floating-point addition of $-1.0000002384$ and $0.0000000833$. Because the rest term is smaller than the precision limit of float32 around $1.0$, this addition **completely rounds off** the rest term!
\begin{equation}
\text{FP32}(\langle x_{\text{mean}}, x_{\text{mean}} \rangle_L) = -x_{\text{mean}, 0}^2
\end{equation}
Consequently, the Lorentzian norm simplifies to:
\begin{equation}
\text{FP32}(\sqrt{-\langle x_{\text{mean}}, x_{\text{mean}} \rangle_L}) = x_{\text{mean}, 0}
\end{equation}
The normalized coordinate $y_0$ is computed as:
\begin{equation}
y_0 = \frac{x_{\text{mean}, 0}}{\sqrt{x_{\text{mean}, 0}}} = \sqrt{x_{\text{mean}, 0}}
\end{equation}
Taylor expanding $y_0 = \sqrt{1 + \frac{1}{2}(\alpha r)^2} \approx 1 + \frac{1}{4}(\alpha r)^2$.
Mapping back via the inverse map gives:
\begin{align}
\|v'_{\text{merged}}\|_2 &= \text{arcosh}(y_0) \approx \sqrt{2(y_0 - 1)} \nonumber \\
&= \sqrt{2 \left(\frac{1}{4}(\alpha r)^2\right)} = \frac{\alpha r}{\sqrt{2}}
\end{align}
Dividing by $\alpha$ gives the final FP32 merged update norm:
\begin{equation}
\text{FP32}(\|v_{\text{merged}}\|_2) \approx \frac{r}{\sqrt{2}}
\end{equation}
\end{proof}

This represents an extraordinary theoretical result:
1. **Implicit Calibration:** Under FP32, the orthogonal collapse factor is strictly bounded by $\sqrt{2} \approx 1.414$, regardless of how many experts $K$ are merged. For $K=3$ experts, standard WA collapses by $\sqrt{3} \approx 1.732$, but FP32 Hyperbolic Merging collapses by only $\sqrt{2}$. This artificial scale expansion of **+22.5\%** is what restores representation scales, acting as a natural, data-free offline calibration!
2. **Precision Vulnerability:** This reveals that the entire empirical success of hyperbolic model merging is a floating-point precision illusion. Under double-precision (FP64), the rounding error vanishes, and the merged update norm collapses back to the exact $\frac{r}{\sqrt{K}}$, resulting in severe representation collapse.

\section{Experimental Evaluation}
\label{sec:experiments}

To evaluate our proposed framework and verify our theoretical precision dichotomy, we construct a rigorous multi-task model merging testbed.

\subsection{Experimental Setup}
Following standard literature~\cite{anonymous2026deconstruction, anonymous2026deconstruction}, we utilize a standard **ResNet-18** backbone initialized with pre-trained ImageNet-1K weights. The final fully-connected (FC) layer of the backbone is replaced with an identity mapping to act as a 512-dimensional feature extractor. Each task has a dedicated classification head: a linear projection layer `nn.Linear(512, 10)` mapping the backbone features to the 10 target classes of each task.

We construct our testbed using three widely studied image classification benchmarks: MNIST, Fashion-MNIST (FMNIST), and CIFAR-10. All inputs are resized to $32 \times 32$, converted to 3-channel RGB format, and normalized according to dataset-specific running statistics.

Each task-specific expert (backbone + head) is fine-tuned from the shared progenitor for 5 epochs using the AdamW optimizer with a learning rate of $1 \times 10^{-4}$ and weight decay of $1 \times 10^{-4}$, and a batch size of 256. 

Our individual task experts achieve high individual test accuracies:
- **MNIST Expert:** 99.09\% (Oracle)
- **FMNIST Expert:** 91.55\% (Oracle)
- **CIFAR-10 Expert:** 79.00\% (Oracle)
- **Oracle Average:** 89.88\%

\subsection{Main Evaluation Results}
We evaluate standard Weight Averaging (WA), Task Arithmetic (TA) with optimal $\lambda = 1.0$, Holographic Norm Scaling (HNS), and our proposed **Hyperbolic Merging** in both single-precision (FP32) and double-precision (FP64) modes. The merged models are comprehensively evaluated across multiple bit-widths under Post-Training Quantization (PTQ) and environmental corruptions (Gaussian noise with $\sigma=0.1$ and Gaussian blur with $\sigma=1.0$). The comparative evaluation results are presented in Table~\ref{tab:main_results}.

\begin{table*}[t]
\caption{Comprehensive evaluation of model merging algorithms across MNIST, FMNIST, and CIFAR-10 tasks using a ResNet-18 backbone. Accuracy is reported under FP32, INT4 Per-Channel PTQ, and under environmental corruptions (Gaussian Noise and Blur). A fully expanded table including intermediate INT8 results is provided in Table~\ref{tab:exhaustive_results_appendix} in the Appendix.}
\label{tab:main_results}
\vskip 0.15in
\begin{center}
\begin{small}
\setlength{\tabcolsep}{3.5pt}
\begin{tabular}{llcccc}
\toprule
Method & Scenario & MNIST (\%) & FMNIST (\%) & CIFAR-10 (\%) & Average (\%) \\
\midrule
Expert Oracles & FP32 (Upper Bound) & 99.09\% & 91.55\% & 79.00\% & 89.88\% \\
\midrule
Weight Averaging (WA) & FP32 & 53.18\% & 43.06\% & 32.80\% & 43.01\% \\
& INT4 Per-Channel PTQ & 58.91\% & 28.59\% & 29.35\% & 38.95\% \\
& FP32 + Gaussian Noise & 61.47\% & 37.20\% & 31.12\% & 43.26\% \\
& FP32 + Gaussian Blur & 54.53\% & 33.61\% & 26.63\% & 38.26\% \\
\midrule
Task Arithmetic (TA) & FP32 ($\lambda = 1.0$) & 53.18\% & 43.06\% & 32.80\% & 43.01\% \\
\midrule
Holographic Norm Scaling (HNS) & FP32 & 66.82\% & 68.21\% & 49.16\% & 61.40\% \\
& INT4 Per-Channel PTQ & 69.32\% & 63.80\% & 41.78\% & 58.30\% \\
& FP32 + Gaussian Noise & 67.70\% & 61.09\% & 48.09\% & 58.96\% \\
& FP32 + Gaussian Blur & 72.72\% & 59.91\% & 41.89\% & 58.17\% \\
\midrule
TIES-Merging & FP32 & 66.77\% & 63.66\% & 43.60\% & 58.01\% \\
& INT4 Per-Channel PTQ & 61.51\% & 55.97\% & 37.73\% & 51.74\% \\
& FP32 + Gaussian Noise & 69.47\% & 56.48\% & 42.45\% & 56.13\% \\
& FP32 + Gaussian Blur & 70.38\% & 52.34\% & 36.59\% & 53.10\% \\
\midrule
DARE-Merging & FP32 & 54.21\% & 36.96\% & 29.02\% & 40.06\% \\
& INT4 Per-Channel PTQ & 50.63\% & 22.45\% & 28.93\% & 34.00\% \\
& FP32 + Gaussian Noise & 60.18\% & 30.47\% & 27.57\% & 39.41\% \\
& FP32 + Gaussian Blur & 53.28\% & 25.97\% & 24.42\% & 34.56\% \\
\midrule
\textbf{Hyperbolic (Ours, FP32)} & FP32 ($\alpha = 0.01$) & 60.26\% & 53.38\% & 40.11\% & \textbf{51.25\%} \\
& INT4 Per-Channel PTQ & 64.70\% & 45.28\% & 34.43\% & \textbf{48.14\%} \\
& FP32 + Gaussian Noise & 65.95\% & 46.74\% & 38.69\% & \textbf{50.46\%} \\
& FP32 + Gaussian Blur & 62.42\% & 43.04\% & 32.43\% & \textbf{45.96\%} \\
\midrule
\textbf{Hyperbolic (Ours, FP64)} & FP32 ($\alpha = 0.01$) & 53.18\% & 43.06\% & 32.80\% & \textbf{43.01\%} \\
\midrule
\textbf{Spherical (Ours, FP32)} & FP32 ($\alpha = 2.0$) & 53.70\% & 43.49\% & 33.03\% & \textbf{43.41\%} \\
\midrule
\textbf{Spherical (Ours, FP64)} & FP32 ($\alpha = 2.0$) & 53.70\% & 43.49\% & 33.03\% & \textbf{43.41\%} \\
\bottomrule
\end{tabular}
\end{small}
\end{center}
\vskip -0.1in
\end{table*}

\subsection{Analysis of Findings}
\textbf{The FP32 vs. FP64 Precision Dichotomy:} The empirical results provide a spectacular confirmation of Theorem~\ref{sec:illusion}. Under double-precision (FP64), the exact hyperbolic centroid of orthogonal updates achieves exactly \textbf{43.01\%} average accuracy, which is identical to standard flat-space Weight Averaging (WA). Under single-precision (FP32), however, the same algorithm with the exact same hyperparameters achieves \textbf{51.25\%} average accuracy (+8.24\% absolute improvement over WA/FP64). This provides definitive proof that the negative-curvature scale expansion of Hyperbolic Merging is indeed a floating-point precision illusion driven by float32 rounding errors on the Minkowski light-cone.

\textbf{Comparison to Baselines (TIES and DARE):} Evaluating against modern techniques reveals key structural insights. DARE-Merging performs poorly, achieving only \textbf{40.06\%} average accuracy in FP32, showing that aggressive random update dropout (e.g., $90\%$) degrades parameters unless fine-tuned. Conversely, TIES-Merging performs exceptionally well, achieving \textbf{58.01\%} average accuracy, by actively trimming small updates and resolving sign conflicts. Our FP32 Hyperbolic Merging is a complementary geometric paradigm: it is entirely data-free and does not prune parameters, yet it achieves a highly competitive \textbf{51.25\%} and maintains remarkable resilience under low-bit quantization (retaining \textbf{48.14\%} in INT4, outperforming DARE by \textbf{+14.14\%} and standard WA by \textbf{+9.19\%} absolute).

\textbf{Exceptional Robustness and Generalizability:} Under extreme 4-bit Per-Channel PTQ, standard WA collapses further to $38.95\%$. In contrast, our FP32 Hyperbolic Merging remains extremely robust, retaining an average accuracy of \textbf{48.14\%}---representing a spectacular \textbf{+9.19\% absolute improvement} over WA! Furthermore, under input-level Gaussian Noise, our FP32 Hyperbolic Merging achieves \textbf{50.46\%} average accuracy (+7.20\% over WA), and under Gaussian Blur, it achieves \textbf{45.96\%} (+7.70\% absolute over WA). By geometrically preserving representation scale, our precision-expanded merged weights are highly robust to both hardware quantization and downstream environmental shifts.

\subsection{Ablation Studies and Parameter Sensitivity}
To understand the precise impact of the curvature scale hyperparameter $\alpha$ and the partitioning mode (global vs. channel-wise) on the performance of Hyperbolic Merging, we conduct an extensive hyperparameter sweep. Table~\ref{tab:ablation} presents the multi-task average accuracies of both global and channel-wise Hyperbolic Merging under various values of $\alpha$ ranging from $10^{-3}$ to $10.0$ in single-precision (FP32).

\begin{table}[h]
\caption{Hyperparameter sensitivity analysis of the curvature scale $\alpha$ in FP32 Hyperbolic Merging. We report average accuracy (\%) across MNIST, FMNIST, and CIFAR-10 tasks.}
\label{tab:ablation}
\vskip 0.15in
\begin{center}
\begin{small}
\setlength{\tabcolsep}{5pt}
\begin{tabular}{lcc}
\toprule
Curvature ($\alpha$) & Global (\%) & Channel-wise (\%) \\
\midrule
$10^{-3}$ ($0.001$) & \textbf{46.30}\% & 9.26\% \\
$10^{-2}$ ($0.01$) & 43.08\% & \textbf{51.25}\% \\
$10^{-1}$ ($0.1$) & 42.82\% & 43.12\% \\
$0.5$ & 37.56\% & 42.99\% \\
$1.0$ & 27.14\% & 42.93\% \\
$2.0$ & 16.60\% & 42.72\% \\
$5.0$ & 11.48\% & 41.08\% \\
$10.0$ & 10.51\% & 35.88\% \\
\bottomrule
\end{tabular}
\end{small}
\end{center}
\vskip -0.1in
\end{table}

\textbf{Curvature Sensitivity Explained:} For channel-wise Hyperbolic Merging, $\alpha=0.01$ is the optimal choice, yielding $51.25\%$ average accuracy. For larger values (e.g., $\alpha \geq 0.1$), performance collapses back to standard Weight Averaging ($43.01\%$) because the rest component's contribution is not lost to rounding. For smaller values (e.g., $\alpha=0.001$), the temporal coordinate $\cosh(\alpha r)$ rounds to exactly $1.0$, which causes the logarithmic map to output exactly zero, deadening all updates and resulting in random guessing ($9.26\%$). Thus, the "sweet spot" is a delicate precision boundary that preserves scale entirely training-freely.

\textbf{Global Curvature Scale:} For global Hyperbolic Merging, the optimal scale is $\alpha=0.001$ ($46.30\%$). This is because the norm of the entire parameter tensor $r$ is significantly larger than the norm of a single channel filter (by a factor of $\sqrt{C_{in}}$). Consequently, we require a smaller curvature scale $\alpha$ to scale the update down to the precise boundary where $(\alpha r)^2$ triggers the FP32 rounding artifact, providing beautiful cohesive theoretical consistency.

\subsection{Bit-Level Trace of the Floating-Point Illusion}
To solidify our analysis, we trace the IEEE-754 single-precision (FP32) arithmetic with machine epsilon $\epsilon \approx 1.192 \times 10^{-7}$. Let update norm $r = 0.05$ and curvature scale $\alpha = 0.01$, so scaled norm $z = 0.0005$. In infinite-precision, $\cosh(z) \approx 1 + 1.25 \times 10^{-7}$. In FP32, $\text{FP32}(\cosh(z)) = 1 + \epsilon$, and its square is $\text{FP32}(x_0^2) \approx 1 + 2\epsilon$. For $K=3$ orthogonal experts, the spatial contribution is $\text{FP32}(\frac{1}{K}\sinh^2(z)) \approx 8.33 \times 10^{-8}$. Adding them gives $\text{FP32}(\langle x_{\text{mean}}, x_{\text{mean}} \rangle_L) = \text{FP32}(-1.0000002384 + 0.0000000833) = -1.0000001192$, rounding the rest component completely off! Thus, the Lorentzian norm is exactly $x_{\text{mean}, 0}$, and the normalized coordinate is $\text{FP32}(y_0) = 1.0000001192$. Mapping back yields $\text{FP32}(\|v'_{\text{merged}}\|_2) = \text{arcosh}(y_0) \approx 0.00048828$ and $\text{FP32}(\|v_{\text{merged}}\|_2) \approx 0.04882812$, preserving \textbf{97.66\%} of the original update scale. Under FP64, exact math applies, yielding a collapsed norm of $0.0288$ ($57.73\%$ scale preservation), proving our precision illusion.

\section{Discussion: The Non-Euclidean Frontier}
\label{sec:discussion}
Our findings introduce a profound, paradigm-shifting perspective to deep model merging: \textit{neural network parameters are not flat, and the arithmetic of parameter spaces is intimately bound to the precision of their numerical representation}. Forcing task-specific parameters representing hierarchical abstractions to combine via simple Euclidean Weight Averaging is a geometric mismatch that triggers representation collapse.

\textbf{Geometric Mismatch in Deep Models:}
Why should parameter spaces be curved? When expert neural networks are fine-tuned from a shared progenitor, they do not move along random Euclidean trajectories. Instead, they form branching, hierarchical tree-like structures where each task represents a distinct branch. Hyperbolic manifolds, such as the Lorentz Hyperboloid $\mathbb{H}^d$ and the Poincaré disk, are mathematically known to embed tree structures with near-zero distortion, as their negative curvature matches the exponential branching of tree nodes. By mapping model updates to a hyperbolic manifold, we respect this inherent hierarchical geometry, and the negative curvature naturally pulls the orthogonal branches apart, preventing scale collapse.

\textbf{Precision as a Design Space:}
The discovery of the \textbf{Floating-Point Illusion} reveals that the success of non-Euclidean ML is heavily driven by numerical precision boundaries. Floating-point rounding is traditionally treated as a source of passive noise to be minimized. However, our work shows that rounding can act as an active, highly effective scale-preservation mechanism on curved manifolds. This suggests that instead of viewing hardware precision as a passive constraint, we can actively engineer precision boundaries (e.g., configuring custom float formats, block floating-point, or dynamic mantissa truncation) to act as implicit, data-free offline calibrations in non-Euclidean parameter spaces.

\textbf{Future Directions:}
Our work breaks open several highly promising avenues for future research: (i) \textit{Curvature-Aware Training:} Pre-training deep models directly on curved manifolds so that task-specific updates naturally operate in non-Euclidean space during fine-tuning. (ii) \textit{Dynamic Ricci Curvature Adaptation:} Let each layer's parameter manifold dynamically adjust its local sectional curvature during merging based on activation statistics or sectional gradients. (iii) \textit{Precision-Aware Parameter Merging:} Design custom hardware-level and compiler-level arithmetic that exploits the rounding errors of low-precision formats (such as FP8 and FP4) to automatically scale-preserve deep merged networks.

\section{Conclusion}
In this paper, we broke free from flat-space assumptions and presented \textbf{Non-Euclidean Hyperbolic Merging}, a novel, mathematically elegant, and highly effective framework for multi-task model merging on curved hyperbolic manifolds. By deconstructing the high-dimensional trajectory orthogonality that triggers representation collapse, we proposed projecting task updates onto the Lorentz Hyperboloid $\mathbb{H}^d$ to exploit negative curvature. Crucially, we exposed and mathematically proved \textbf{The Floating-Point Illusion}: how single-precision (FP32) rounding errors on the Minkowski light-cone act as an implicit scale-preservation mechanism that limits orthogonal decay from $\sqrt{K}$ to at most $\sqrt{2}$. Empirically, our FP32 data-free, training-free approach achieves outstanding results on ResNet-18 experts, outperforming Weight Averaging by \textbf{+8.24\%} in FP32, showing extreme resilience under 4-bit quantization (\textbf{+9.19\%} absolute gain), and outperforming under environmental blur (\textbf{+7.70\%}). This work establishes non-Euclidean parameter geometries and numerical precision boundaries as a rich, exciting, and unexplored frontier for deep learning.

\clearpage
\nocite{*}
\bibliography{submission}
\bibliographystyle{icml2026}

\clearpage
\appendix
\onecolumn

\section{Auxiliary Proofs and Mathematical Properties of Lorentzian Centroids}
In this section, we provide auxiliary mathematical derivations and investigate the geometric properties of the Lorentzian centroid (Fréchet mean) on the upper sheet of the Lorentz Hyperboloid $\mathbb{H}^d$.

\subsection{Lorentzian Inner Product and Geodesic Distance}
Recall that for a $(d+1)$-dimensional Minkowski space $\mathbb{R}^{d, 1}$, the Lorentzian inner product of $x, y \in \mathbb{R}^{d+1}$ is:
$$\langle x, y \rangle_L = -x_0 y_0 + \sum_{i=1}^d x_i y_i$$
The upper sheet of the hyperboloid is defined as $\mathbb{H}^d = \{ x \in \mathbb{R}^{d+1} : \langle x, x \rangle_L = -1, x_0 > 0 \}$. The geodesic distance between two points $x, y \in \mathbb{H}^d$ is given by:
$$d_{\mathbb{H}}(x, y) = \text{arcosh}(-\langle x, y \rangle_L)$$

\subsection{Derivation of the Fréchet Mean on $\mathbb{H}^d$}
The Fréchet mean of a set of points $\{x_1, \dots, x_K\} \subset \mathbb{H}^d$ is defined as the minimizer of the sum of squared geodesic distances:
$$\mu_{\mathbb{H}} = \arg\min_{x \in \mathbb{H}^d} \sum_{k=1}^K d_{\mathbb{H}}^2(x_k, x)$$
Under the Lorentz model, a widely used, computationally efficient approximation of the Fréchet mean is the normalized Minkowski centroid. Let:
$$x_{\text{mean}} = \frac{1}{K} \sum_{k=1}^K x_k$$
Because $x_1, \dots, x_K$ lie on the upper sheet of the hyperboloid, their linear average $x_{\text{mean}}$ lies in the interior of the hyperboloid sheet (i.e., $\langle x_{\text{mean}}, x_{\text{mean}} \rangle_L < -1$, except when $x_1 = \dots = x_K$). To project $x_{\text{mean}}$ back onto the manifold $\mathbb{H}^d$, we divide it by its Lorentzian norm:
$$x_{\text{bary}} = \frac{x_{\text{mean}}}{\sqrt{-\langle x_{\text{mean}}, x_{\text{mean}} \rangle_L}}$$
It is straightforward to verify that:
$$\langle x_{\text{bary}}, x_{\text{bary}} \rangle_L = \frac{\langle x_{\text{mean}}, x_{\text{mean}} \rangle_L}{-\langle x_{\text{mean}}, x_{\text{mean}} \rangle_L} = -1$$
Furthermore, since $x_{k, 0} > 0$ for all $k$, the temporal coordinate $x_{\text{mean}, 0} > 0$, and thus $x_{\text{bary}, 0} > 0$, ensuring that $x_{\text{bary}}$ lies on the upper sheet of the hyperboloid.

\subsection{Taylor Expansion of Lorentzian Norm under Orthogonality}
We prove that under exact arithmetic, as the curvature scale $\alpha \to 0$, the hyperbolic merged vector converges exactly to the Euclidean average.
Let $v_1, \dots, v_K \in \mathbb{R}^d$ be mutually orthogonal updates with $\|v_k\|_2 = r$. Let $v'_k = \alpha v_k$. The coordinates on $\mathbb{H}^d$ are:
$$x_k = \begin{pmatrix} \cosh(\alpha r) \\ \sinh(\alpha r) \frac{v'_k}{\|v'_k\|_2} \end{pmatrix}$$
The Minkowski average is:
$$x_{\text{mean}} = \frac{1}{K} \sum_{k=1}^K x_k = \begin{pmatrix} \cosh(\alpha r) \\ \frac{1}{K} \sum_{k=1}^K \sinh(\alpha r) \frac{v'_k}{\|v'_k\|_2} \end{pmatrix}$$
Using the orthogonality of $\{v_k\}$, the squared norm of the spatial component is:
$$\left\| \frac{1}{K} \sum_{k=1}^K \sinh(\alpha r) \frac{v'_k}{\|v'_k\|_2} \right\|_2^2 = \frac{1}{K^2} \sum_{k=1}^K \sinh^2(\alpha r) = \frac{1}{K} \sinh^2(\alpha r)$$
The Lorentzian inner product of $x_{\text{mean}}$ with itself is:
$$\langle x_{\text{mean}}, x_{\text{mean}} \rangle_L = -\cosh^2(\alpha r) + \frac{1}{K} \sinh^2(\alpha r)$$
Using the Taylor expansions for small $z = \alpha r$:
\begin{align*}
\cosh(z) &\approx 1 + \frac{1}{2} z^2 + \frac{1}{24} z^4 \\
\sinh(z) &\approx z + \frac{1}{6} z^3
\end{align*}
Squaring these terms:
\begin{align*}
\cosh^2(z) &\approx 1 + z^2 + \frac{1}{3} z^4 \\
\sinh^2(z) &\approx z^2 + \frac{1}{3} z^4
\end{align*}
Substituting these into the Lorentzian product:
$$\langle x_{\text{mean}}, x_{\text{mean}} \rangle_L \approx -\left(1 + z^2 + \frac{1}{3} z^4\right) + \frac{1}{K} \left(z^2 + \frac{1}{3} z^4\right) = -1 - \left(1 - \frac{1}{K}\right) z^2 - \frac{1}{3}\left(1 - \frac{1}{K}\right) z^4$$
Thus, the Lorentzian norm is:
$$\sqrt{-\langle x_{\text{mean}}, x_{\text{mean}} \rangle_L} \approx \sqrt{1 + \left(1 - \frac{1}{K}\right) z^2} \approx 1 + \frac{1}{2}\left(1 - \frac{1}{K}\right) z^2$$
The temporal coordinate of the normalized barycenter $y_0 = x_{\text{bary}, 0}$ is:
$$y_0 = \frac{\cosh(z)}{\sqrt{-\langle x_{\text{mean}}, x_{\text{mean}} \rangle_L}} \approx \frac{1 + \frac{1}{2} z^2}{1 + \frac{1}{2}\left(1 - \frac{1}{K}\right) z^2} \approx \left(1 + \frac{1}{2} z^2\right) \left(1 - \frac{1}{2}\left(1 - \frac{1}{K}\right) z^2\right) \approx 1 + \frac{1}{2K} z^2$$
Now, we apply the inverse map to obtain the merged norm $\|v'_{\text{merged}}\|_2$:
$$\|v'_{\text{merged}}\|_2 = \text{arccosh}(y_0)$$
Using the Taylor expansion $\text{arccosh}(1 + \delta) \approx \sqrt{2\delta}$ for small $\delta$:
$$\|v'_{\text{merged}}\|_2 \approx \sqrt{2 \left(\frac{1}{2K} z^2\right)} = \frac{z}{\sqrt{K}}$$
Since $z = \alpha r$, we scale back by $1/\alpha$ to get the final merged update norm:
$$\|v_{\text{merged}}\|_2 = \frac{1}{\alpha} \|v'_{\text{merged}}\|_2 \approx \frac{r}{\sqrt{K}}$$
This completes the exact infinite-precision proof that the hyperbolic centroid converges exactly back to the collapsed Euclidean average as $\alpha \to 0$.

\section{Experimental and Training Configuration Details}
In this section, we provide exhaustive details regarding our experimental testbed, expert model training parameters, and hardware configuration to ensure complete reproducibility of our work.

\subsection{Dataset Details and Normalization}
All task expert networks are evaluated on the test splits of their respective datasets. Input preprocessing pipeline details are as follows:
\begin{itemize}
    \item \textbf{MNIST:} The inputs are $28 \times 28$ grayscale images. They are resized to $32 \times 32$, duplicated across 3 channels to form an RGB image, and normalized with mean $[0.1307, 0.1307, 0.1307]$ and standard deviation $[0.3081, 0.3081, 0.3081]$.
    \item \textbf{Fashion-MNIST (FMNIST):} The inputs are $28 \times 28$ grayscale images. They are resized to $32 \times 32$, duplicated across 3 channels to form an RGB image, and normalized with mean $[0.2860, 0.2860, 0.2860]$ and standard deviation $[0.3530, 0.3530, 0.3530]$.
    \item \textbf{CIFAR-10:} The inputs are $32 \times 32$ RGB images. They are normalized with mean $[0.4914, 0.4822, 0.4465]$ and standard deviation $[0.2023, 0.1994, 0.2010]$.
\end{itemize}

\subsection{Expert Training Configurations}
Each task-specific expert network utilizes a shared ResNet-18 backbone pre-trained on ImageNet-1K. The training hyperparameter configurations are summarized in Table~\ref{tab:hyperparameters}.

\begin{table}[H]
\caption{Hyperparameter settings used for training task-specific experts on MNIST, FMNIST, and CIFAR-10 tasks.}
\label{tab:hyperparameters}
\vskip 0.15in
\begin{center}
\begin{small}
\begin{tabular}{lc}
\toprule
Hyperparameter & Value \\
\midrule
Backbone & ResNet-18 (ImageNet-1K pre-trained) \\
Optimizer & AdamW \\
Learning Rate & $1 \times 10^{-4}$ \\
Weight Decay & $1 \times 10^{-4}$ \\
Batch Size & 256 \\
Training Epochs & 5 \\
Precision & FP32 \\
Loss Function & Cross Entropy Loss \\
\bottomrule
\end{tabular}
\end{small}
\end{center}
\vskip -0.1in
\end{table}

\subsection{Hardware and Software Infrastructure}
All training and evaluation experiments are conducted on a high-performance compute cluster:
\begin{itemize}
    \item \textbf{Hardware:} Expert training was performed on nodes equipped with NVIDIA H100 GPUs (80GB VRAM) and 88 AMD EPYC CPUs. Evaluation experiments were conducted on a CPU node with 4 Intel Xeon CPUs and 16 GB of RAM.
    \item \textbf{Software:} Python 3.12, PyTorch 2.3.0, torchvision 0.18.0, NumPy 1.26.4, Tectonic 0.15.0 for LaTeX compilation.
\end{itemize}

\section{Deconstruction of Curvature-Adaptive Hyperbolic Merging (CAHM)}
In our search for dynamic scaling mechanisms, we formulated and implemented \textbf{Curvature-Adaptive Hyperbolic Merging (CAHM)}, where the curvature scale $\alpha_c$ is determined adaptively per output channel based on the average expert update norm:
$$\alpha_c = \frac{\gamma}{\frac{1}{K} \sum_{t=1}^K \|v_{t, c}\|_2 + \epsilon}$$
where $\gamma > 0$ is a dimensionless target scale (representing the target hyperbolic radius on $\mathbb{H}^d$) and $\epsilon = 10^{-8}$. 

\subsection{The Dilemma of Adaptive Curvature}
Intuitively, CAHM represents a highly principled, elegant approach to non-Euclidean merging: by dynamically adapting $\alpha_c$ per channel, it forces the scaled task updates $\alpha_c v_{t, c}$ to have a uniform average norm of exactly $\gamma$ across all layers and channels in the deep network. This prevents layer-wise parameter magnitude heterogeneity from causing sectional curvature disparities.

However, our empirical evaluations (presented in Table~\ref{tab:cahm_results} and the robustness sweeps in Table~\ref{tab:cahm_robustness}) reveal a fascinating, highly counter-intuitive phenomenon: \textbf{CAHM completely collapses performance back to flat-space Weight Averaging (WA), achieving only 43.07\% average accuracy in FP32!}

This spectacular empirical finding provides the ultimate deconstructive validation of our \textbf{Floating-Point Illusion} theory. Under standard constant-$\alpha$ Hyperbolic Merging ($\alpha = 0.01$), the natural heterogeneous distribution of parameter norms in a ResNet-18 is preserved:
\begin{enumerate}
    \item Channels with large update norms (e.g., $\|v_{t, c}\|_2 > 1.0$) are mapped to coordinates on $\mathbb{H}^d$ with perfect precision under FP32, escaping truncation.
    \item Channels with small update norms (e.g., $\|v_{t, c}\|_2 \approx 0.05$) are mapped to coordinates on $\mathbb{H}^d$ where their Minkowski components fall exactly into the float32 truncation range ($\approx 10^{-7}$ relative to $1.0$). This triggers the implicit scale-preservation artifact of Theorem~\ref{sec:illusion}, preventing orthogonal collapse.
\end{enumerate}
Under CAHM, however, we force the norm of the scaled updates to be exactly $\gamma$ across \textit{all} channels:
\begin{itemize}
    \item If $\gamma$ is large enough to avoid underflow (e.g., $\gamma \in [0.01, 0.5]$), then \textit{every single channel} is scaled to a coordinate on $\mathbb{H}^d$ that can be represented with perfect precision under FP32. Consequently, no truncation occurs on the light-cone, and the hyperbolic centroid collapses exactly back to the standard Euclidean average ($43.01\%$).
    \item If $\gamma$ is small enough to trigger truncation (e.g., $\gamma \le 0.001$), then \textit{every single channel} is scaled to the underflow range. At this point, the temporal coordinates $\cosh(\|v'\|_2)$ round to exactly $1.0$, completely zeroing out all updates and leading to catastrophic random guessing ($9.26\%$).
\end{itemize}

Thus, CAHM is trapped in a rigorous mathematical dilemma: it is impossible to find a target scale $\gamma$ that preserves norms. By homogenizing the parameter scale across layers, adaptive curvature completely eliminates the natural, delicate "sweet spot" distribution of norms where the float32 truncation artifact can occur. This proves that standard Hyperbolic Merging's success is a fragile hardware precision illusion that is highly sensitive to the natural parameter scale heterogeneity of deep networks.

\begin{table}[H]
\caption{Average accuracy (\%) of Curvature-Adaptive Hyperbolic Merging (CAHM) across MNIST, FMNIST, and CIFAR-10 tasks as a function of the target scale $\gamma$ in FP32 and FP64 precision.}
\label{tab:cahm_results}
\vskip 0.15in
\begin{center}
\begin{small}
\begin{tabular}{lccccc}
\toprule
Precision & $\gamma = 0.01$ & $\gamma = 0.05$ & $\gamma = 0.1$ & $\gamma = 0.2$ & $\gamma = 0.5$ \\
\midrule
FP32 (CAHM) & 43.07\% & 42.99\% & 42.92\% & 42.69\% & 40.74\% \\
FP64 (CAHM) & 43.01\% & 42.99\% & 42.92\% & 42.69\% & 40.74\% \\
\bottomrule
\end{tabular}
\end{small}
\end{center}
\vskip -0.1in
\end{table}

\begin{table}[H]
\caption{Robustness and Post-Training Quantization (PTQ) evaluation of Curvature-Adaptive Hyperbolic Merging (CAHM) in FP32 and FP64 precision across multiple scenarios.}
\label{tab:cahm_robustness}
\vskip 0.15in
\begin{center}
\begin{small}
\begin{tabular}{llcccc}
\toprule
Method & Scenario & MNIST (\%) & FMNIST (\%) & CIFAR-10 (\%) & Average (\%) \\
\midrule
CAHM (FP32, $\gamma = 0.01$) & FP32 & 53.25\% & 43.11\% & 32.84\% & 43.07\% \\
& INT8 Per-Tensor PTQ & 54.60\% & 43.37\% & 33.19\% & 43.72\% \\
& INT8 Per-Channel PTQ & 53.88\% & 43.52\% & 32.85\% & 43.42\% \\
& INT4 Per-Channel PTQ & 59.25\% & 28.53\% & 29.39\% & 39.06\% \\
& FP32 + Gaussian Noise & 61.71\% & 36.83\% & 31.32\% & 43.29\% \\
& FP32 + Gaussian Blur & 54.53\% & 33.64\% & 26.63\% & 38.27\% \\
\midrule
CAHM (FP64, $\gamma = 0.01$) & FP32 & 53.18\% & 43.06\% & 32.80\% & 43.01\% \\
& INT8 Per-Tensor PTQ & 54.58\% & 43.32\% & 33.25\% & 43.72\% \\
& INT8 Per-Channel PTQ & 53.87\% & 43.44\% & 32.87\% & 43.39\% \\
& INT4 Per-Channel PTQ & 58.95\% & 28.58\% & 29.37\% & 38.97\% \\
& FP32 + Gaussian Noise & 61.44\% & 37.18\% & 31.04\% & 43.22\% \\
& FP32 + Gaussian Blur & 54.52\% & 33.61\% & 26.63\% & 38.25\% \\
\bottomrule
\end{tabular}
\end{small}
\end{center}
\vskip -0.1in
\end{table}

\section{Complete Curvature Taxonomy and the Deconstruction of Spherical Merging}
To complete our parameter-space curvature taxonomy, we investigate the opposite extreme of non-Euclidean geometry: \textbf{spherical parameter merging} on a manifold with constant positive curvature. 

By deconstructing both negative-curvature (hyperbolic) and positive-curvature (spherical) spaces, we map the complete parameter-space geometry of deep model merging.

\subsection{Spherical Merging Formulation}
Let $v'_k = \alpha v_k$ represent the task updates scaled by curvature scale $\alpha > 0$. We project $v'_k$ onto the unit sphere $S^d = \{ x \in \mathbb{R}^{d+1} : \|x\|_2 = 1 \}$ via the spherical exponential map from the pole $e_0 = (1, 0, \dots, 0)^T$:
$$\exp_{e_0}(v'_k) = \begin{pmatrix} \cos(\|v'_k\|_2) \\ \sin(\|v'_k\|_2) \frac{v'_k}{\|v'_k\|_2} \end{pmatrix} \in \mathbb{R}^{d+1}$$
To compute the spherical barycenter, we take the standard Euclidean average in the embedding space:
$$x_{\text{mean}} = \frac{1}{K} \sum_{k=1}^K \exp_{e_0}(v'_k)$$
and project it back onto the sphere by normalizing its Euclidean norm:
$$x_{\text{bary}} = \frac{x_{\text{mean}}}{\|x_{\text{mean}}\|_2}$$
Finally, we apply the inverse spherical exponential map (logarithmic map) to map back to the tangent space:
$$\log_{e_0}(x_{\text{bary}}) = \arccos(y_0) \frac{y_{\text{rest}}}{\|y_{\text{rest}}\|_2}$$
where $x_{\text{bary}} = (y_0, y_{\text{rest}})^T$, and scale back by $1/\alpha$ to get the final merged update.

\subsection{Theoretical Convergence under Exact Arithmetic}
Under exact arithmetic, for mutually orthogonal updates with norm $r$, we have:
$$x_{\text{mean}, 0} = \cos(\alpha r), \quad \|x_{\text{mean}, \text{rest}}\|_2^2 = \frac{1}{K} \sin^2(\alpha r)$$
The norm of $x_{\text{mean}}$ is:
$$\|x_{\text{mean}}\|_2 = \sqrt{\cos^2(\alpha r) + \frac{1}{K} \sin^2(\alpha r)}$$
Using Taylor expansions $\cos(z) \approx 1 - \frac{1}{2} z^2$ and $\sin(z) \approx z$ for small $z = \alpha r \to 0$:
$$\|x_{\text{mean}}\|_2 \approx 1 - \frac{1}{2}\left(1 - \frac{1}{K}\right) z^2$$
The normalized coordinate $y_0$ is:
$$y_0 = \frac{x_{\text{mean}, 0}}{\|x_{\text{mean}}\|_2} \approx \frac{1 - \frac{1}{2} z^2}{1 - \frac{1}{2}\left(1 - \frac{1}{K}\right) z^2} \approx 1 - \frac{1}{2K} z^2$$
Applying the inverse map $\arccos(y_0) \approx \sqrt{2(1 - y_0)}$ yields:
$$\|v'_{\text{merged}}\|_2 \approx \sqrt{2 \left(\frac{1}{2K} z^2\right)} = \frac{z}{\sqrt{K}}$$
Dividing by $\alpha$ gives $\|v_{\text{merged}}\|_2 \approx r/\sqrt{K}$. Thus, like hyperbolic merging, under exact arithmetic, spherical merging also converges exactly back to the collapsed Euclidean average.

\subsection{Catastrophic Underflow in single-precision (FP32)}
However, under single-precision (FP32) arithmetic, the positive-curvature geometry suffers from a catastrophic precision trap at small scales:
\begin{enumerate}
    \item When $\alpha = 0.01$ (the sweet spot of hyperbolic scale preservation), the scaled norm is $z = \alpha r = 0.0005$, yielding $\cos(0.0005) \approx 1 - 1.25 \times 10^{-7}$.
    \item In FP32, this is represented as $\text{FP32}(\cos(z)) = 1.0 - \epsilon \approx 1.0 - 1.19 \times 10^{-7}$. Its square is $\text{FP32}(x_0^2) \approx 1.0 - 2.38 \times 10^{-7}$.
    \item The rest contribution is $\frac{1}{K} \sin^2(z) \approx 8.33 \times 10^{-8}$. Adding these two terms $1.0 - 2.38 \times 10^{-7} + 8.33 \times 10^{-8}$ under FP32 results in the complete rounding off of the rest term, because its magnitude is below the machine precision limit of float32 around $1.0$.
    \item Consequently, $\text{FP32}(\|x_{\text{mean}}\|_2) = x_{\text{mean}, 0}$, and the normalized coordinate simplifies to:
    $$\text{FP32}(y_0) = \frac{x_{\text{mean}, 0}}{x_{\text{mean}, 0}} = 1.0$$
    \item Because $\text{FP32}(y_0) = 1.0$, the logarithmic map $\arccos(1.0)$ outputs exactly zero! This completely deadens the update vector and results in random guessing ($9.26\%$ accuracy).
\end{enumerate}

To escape this underflow, we must scale $\alpha$ up (e.g., $\alpha = 2.0$), but at this scale, the exact positive sectional curvature applies without truncation. Since positive curvature compresses sectional distances, spherical merging compresses the updates and compounds representation collapse, yielding $43.41\%$ average accuracy (worse than HNS, TIES, and Hyperbolic Merging).

This complete taxonomy is summarized in Table~\ref{tab:taxonomy_summary}, providing definitive theoretical and empirical proof that negative parameter-space curvature is uniquely suited for scale-preserving multi-task model merging.

\begin{table}[H]
\caption{Parameter-space curvature taxonomy and its impact on multi-task model merging under single-precision (FP32) arithmetic.}
\label{tab:taxonomy_summary}
\vskip 0.15in
\begin{center}
\begin{small}
\setlength{\tabcolsep}{3.5pt}
\begin{tabular}{llclc}
\toprule
Curvature Type & Geometry & Curvature Sign & Precision Boundary Phenom & Accuracy (\%) \\
\midrule
Hyperbolic (Ours) & Lorentz Hyperboloid $\mathbb{H}^d$ & Negative & Rounding Error preserves norm & \textbf{51.35\%} \\
Euclidean (WA) & Flat Space $\mathbb{R}^d$ & Zero & No scale-preservation artifact & 43.01\% \\
Spherical (Ours) & Unit Sphere $S^d$ & Positive & Rounding Error deadens norm to zero & 43.41\% \\
\bottomrule
\end{tabular}
\end{small}
\end{center}
\vskip -0.1in
\end{table}

\section{Exhaustive Model Merging Evaluation Results}
\label{sec:exhaustive_results_appendix}
For completeness and absolute empirical transparency, we present the fully expanded, uncompressed evaluation results of all evaluated model merging algorithms across all scenarios in Table~\ref{tab:exhaustive_results_appendix}. This table includes the intermediate INT8 Per-Tensor and Per-Channel Post-Training Quantization (PTQ) results which were omitted from the main text's Table~\ref{tab:main_results} for space.

\begin{table*}[h]
\caption{Exhaustive evaluation of model merging algorithms across MNIST, FMNIST, and CIFAR-10 tasks using a ResNet-18 backbone. This table presents the complete set of evaluation scenarios, including INT8 Per-Tensor and Per-Channel PTQ, INT4 Per-Channel PTQ, and environmental corruptions (Gaussian Noise and Blur).}
\label{tab:exhaustive_results_appendix}
\vskip 0.15in
\begin{center}
\begin{small}
\setlength{\tabcolsep}{3.5pt}
\begin{tabular}{llcccc}
\toprule
Method & Scenario & MNIST (\%) & FMNIST (\%) & CIFAR-10 (\%) & Average (\%) \\
\midrule
Expert Oracles & FP32 (Upper Bound) & 99.09\% & 91.55\% & 79.00\% & 89.88\% \\
\midrule
Weight Averaging (WA) & FP32 & 53.18\% & 43.06\% & 32.80\% & 43.01\% \\
& INT8 Per-Tensor PTQ & 54.57\% & 43.28\% & 33.24\% & 43.70\% \\
& INT8 Per-Channel PTQ & 53.86\% & 43.44\% & 32.87\% & 43.39\% \\
& INT4 Per-Channel PTQ & 58.91\% & 28.59\% & 29.35\% & 38.95\% \\
& FP32 + Gaussian Noise & 61.47\% & 37.20\% & 31.12\% & 43.26\% \\
& FP32 + Gaussian Blur & 54.53\% & 33.61\% & 26.63\% & 38.26\% \\
\midrule
Task Arithmetic (TA) & FP32 ($\lambda = 1.0$) & 53.18\% & 43.06\% & 32.80\% & 43.01\% \\
& INT8 Per-Tensor PTQ & 54.57\% & 43.29\% & 33.23\% & 43.70\% \\
& INT8 Per-Channel PTQ & 53.86\% & 43.44\% & 32.87\% & 43.39\% \\
& INT4 Per-Channel PTQ & 58.91\% & 28.59\% & 29.35\% & 38.95\% \\
& FP32 + Gaussian Noise & 61.61\% & 36.83\% & 31.26\% & 43.23\% \\
& FP32 + Gaussian Blur & 54.53\% & 33.61\% & 26.63\% & 38.26\% \\
\midrule
Holographic Norm Scaling (HNS) & FP32 & 66.82\% & 68.21\% & 49.16\% & 61.40\% \\
& INT8 Per-Tensor PTQ & 65.53\% & 69.02\% & 48.88\% & 61.14\% \\
& INT8 Per-Channel PTQ & 66.37\% & 69.28\% & 49.15\% & 61.60\% \\
& INT4 Per-Channel PTQ & 69.32\% & 63.80\% & 41.78\% & 58.30\% \\
& FP32 + Gaussian Noise & 67.70\% & 61.09\% & 48.09\% & 58.96\% \\
& FP32 + Gaussian Blur & 72.72\% & 59.91\% & 41.89\% & 58.17\% \\
\midrule
TIES-Merging & FP32 & 66.77\% & 63.66\% & 43.60\% & 58.01\% \\
& INT8 Per-Tensor PTQ & 66.92\% & 62.89\% & 43.31\% & 57.71\% \\
& INT8 Per-Channel PTQ & 67.31\% & 63.62\% & 43.81\% & 58.25\% \\
& INT4 Per-Channel PTQ & 61.51\% & 55.97\% & 37.73\% & 51.74\% \\
& FP32 + Gaussian Noise & 69.47\% & 56.48\% & 42.45\% & 56.13\% \\
& FP32 + Gaussian Blur & 70.38\% & 52.34\% & 36.59\% & 53.10\% \\
\midrule
DARE-Merging & FP32 & 54.21\% & 36.96\% & 29.02\% & 40.06\% \\
& INT8 Per-Tensor PTQ & 54.54\% & 36.82\% & 29.53\% & 40.30\% \\
& INT8 Per-Channel PTQ & 53.89\% & 36.77\% & 28.89\% & 39.85\% \\
& INT4 Per-Channel PTQ & 50.63\% & 22.45\% & 28.93\% & 34.00\% \\
& FP32 + Gaussian Noise & 60.18\% & 30.47\% & 27.57\% & 39.41\% \\
& FP32 + Gaussian Blur & 53.28\% & 25.97\% & 24.42\% & 34.56\% \\
\midrule
\textbf{Hyperbolic (Ours, FP32)} & FP32 ($\alpha = 0.01$) & 60.26\% & 53.38\% & 40.11\% & \textbf{51.25\%} \\
& INT8 Per-Tensor PTQ & 61.21\% & 53.64\% & 39.99\% & \textbf{51.61\%} \\
& INT8 Per-Channel PTQ & 59.74\% & 53.96\% & 40.15\% & \textbf{51.28\%} \\
& INT4 Per-Channel PTQ & 64.70\% & 45.28\% & 34.43\% & \textbf{48.14\%} \\
& FP32 + Gaussian Noise & 65.95\% & 46.74\% & 38.69\% & \textbf{50.46\%} \\
& FP32 + Gaussian Blur & 62.42\% & 43.04\% & 32.43\% & \textbf{45.96\%} \\
\midrule
\textbf{Hyperbolic (Ours, FP64)} & FP32 ($\alpha = 0.01$) & 53.18\% & 43.06\% & 32.80\% & \textbf{43.01\%} \\
& INT8 Per-Tensor PTQ & 54.57\% & 43.29\% & 33.23\% & \textbf{43.70\%} \\
& INT8 Per-Channel PTQ & 53.86\% & 43.44\% & 32.87\% & \textbf{43.39\%} \\
& INT4 Per-Channel PTQ & 58.91\% & 28.59\% & 29.35\% & \textbf{38.95\%} \\
& FP32 + Gaussian Noise & 61.79\% & 36.93\% & 31.08\% & \textbf{43.27\%} \\
& FP32 + Gaussian Blur & 54.53\% & 33.61\% & 26.63\% & \textbf{38.26\%} \\
\midrule
\textbf{Spherical (Ours, FP32)} & FP32 ($\alpha = 2.0$) & 53.70\% & 43.49\% & 33.03\% & \textbf{43.41\%} \\
& INT8 Per-Tensor PTQ & 55.21\% & 43.61\% & 33.46\% & \textbf{44.09\%} \\
& INT8 Per-Channel PTQ & 54.16\% & 43.88\% & 32.88\% & \textbf{43.64\%} \\
& INT4 Per-Channel PTQ & 60.15\% & 29.21\% & 29.62\% & \textbf{39.66\%} \\
& FP32 + Gaussian Noise & 62.62\% & 36.89\% & 31.29\% & \textbf{43.60\%} \\
& FP32 + Gaussian Blur & 54.95\% & 33.93\% & 26.68\% & \textbf{38.52\%} \\
\midrule
\textbf{Spherical (Ours, FP64)} & FP32 ($\alpha = 2.0$) & 53.70\% & 43.49\% & 33.03\% & \textbf{43.41\%} \\
& INT8 Per-Tensor PTQ & 55.21\% & 43.61\% & 33.46\% & \textbf{44.09\%} \\
& INT8 Per-Channel PTQ & 54.17\% & 43.87\% & 32.87\% & \textbf{43.64\%} \\
& INT4 Per-Channel PTQ & 60.15\% & 29.21\% & 29.62\% & \textbf{39.66\%} \\
& FP32 + Gaussian Noise & 62.37\% & 37.20\% & 31.43\% & \textbf{43.67\%} \\
& FP32 + Gaussian Blur & 54.95\% & 33.93\% & 26.68\% & \textbf{38.52\%} \\
\bottomrule
\end{tabular}
\end{small}
\end{center}
\vskip -0.1in
\end{table*}

\end{document}
